\subsection{模型的主要优点}

第一，四问共享同一条集合证据链。示向误差被保留为半平面可行域，清除条件由最小包围圆覆盖整个可行域，而非把测得方位当作精确射线。问题三、四又用有限覆盖点集保证发现，因此“发现、定位、清除”三个环节分别有可检查的几何条件。

第二，完整性和效率改动相互分离。v2 只改变 61 个必要覆盖点的访问顺序，v3 只使用题面给出的 16 源上界停止；二者都不删除覆盖点、不缩小示向区域，也不把无信号当作定向源不存在。消融实验因此能把路线续接与停止规则的收益分别归因。

第三，模型与通信状态机一致。程序按请求编号保存动作；服务器已经执行而响应丢失时，重试复用相同编号和内容，避免重复移动或清除。退出只发生在 61 点覆盖完成、清除数达到 16 或预算保护触发三类显式状态。六轮官方演练中全部请求被接受，为协议实现提供了独立于合成模拟的检查。

\subsection{模型的局限}

本文的保证性策略较保守。定向源无信号可能来自距离或背向，程序完全不利用它收缩位置集合，因而避免误删真值，但可能多走路、多测量。源圆域用 32 边外接多边形参与计算，也会轻微扩大初始定位区域；这不会破坏包含真源的保证，却可能延迟最小包围圆降到清除阈值。

最近未完成点是局部贪心规则，没有联合优化“覆盖路线、发现后定位绕行、剩余频道数”。表\ref{tab:ablation}证明它优于固定续扫，但没有证明全局最短。20 米网格兜底追求有限完备性，在不利几何下可能产生较多清除尝试，现实运行受 20 分钟和虚拟 100 小时双重预算约束。

合成模拟还依赖脚本给定的源位置、接收半径、方位误差和定向比例分布。这些分布不是题面承诺，因而 100/100 全部清除只能作为回归证据。v3 的“16 源立即退出”已有单元测试和同种子合成对照；现有官方 16 源演练使用的是 v2，尚未取得 v3 在官方 16 源案例上的实机样本。正式测试仍为 0 次，本文不能报告正式成绩或声称当前策略最优。

\subsection{可继续实施的改进}

在不削弱覆盖证明的前提下，可把 61 个点视为必须访问的节点，以当前机器狗位置、未清除频道数和已知定位绕行为代价，使用局部交换或滚动窗口重排最近若干节点。每次优化后仍保留全部未完成点，因而发现完备性不变；是否采用应以同种子配对实验的虚拟时间和最坏值共同判断。

问题二和单源定位可改用候选域上的极大极小选点：枚举下一检测位置，对可行域边界采样真源位置和误差端点，最小化更新后最小包围圆半径的最坏值。该方法能直接对应清除目标，但计算量高于当前横向基线启发式，需要先验证每次计算能稳定小于通信动作间隔。

正式测试结束后，首先从原始日志复算式\eqref{eq:time}，再按问题三、四分别填写三轮表格；随后检查合成结论在正式案例中是否仍成立。若正式数据与演练趋势冲突，应保留原始结果并修改解释，不应通过更换合成随机种子追求一致。
